% ========== 封面与基础信息 ==========
\title{ElegantBook测试文档}
\subtitle{模板核心特性}
\author{测试工程师}
\institute{模板测试组}
\date{\today}
\version{v1.0}
\extrainfo{本文件仅用于模板功能测试，无任何学术价值}
\bioinfo{测试批次}{TEST-FULL-20260506}

% ========== 目录与基础设置 ==========
\setcounter{tocdepth}{3} % 目录显示到3级

% ========== 测试用参考文献 ==========
\addbibresource[location=local]{reference.bib}

\begin{document}

\maketitle
\frontmatter
\tableofcontents

% ========== 前言（带Section分段） ==========
\pagenumbering{Roman}
\setcounter{page}{0}
\chapter*{测试前言}
\markboth{测试前言}{}
\section*{文档说明}
\markright{文档说明}
本文档为 ElegantBook 全功能测试用例，覆盖模板所有基础与高级功能，所有内容均为虚构测试数据，无实际意义。
\newpage
\section*{测试范围}
\markright{测试范围}
本次测试包含：封面、目录、字体、多级列表、复杂公式、全部数学环境、交叉引用、图表、旁注、章节摘要、习题、参考文献、附录。

\mainmatter

% ========== 第1章：基础排版与多级列表 ==========
\chapter{基础排版与列表测试}
\begin{introduction}[本章测试项]
\item 中文字体与文本格式
\item 3级无序列表嵌套
\item 3级有序列表嵌套
\item 普通段落与强调文本
\end{introduction}

\section{字体与文本测试}
正常正文测试。\textbf{加粗}、\textit{斜体}、\underline{下划线}、\textcolor{red}{红色文字测试}。

{\songti 宋体测试：\textbf{加粗}、\textit{斜体}、\underline{下划线}}；\par
{\heiti 黑体测试：\textbf{加粗}、\textit{斜体}、\underline{下划线}}；\par
{\kaishu 楷书测试：\textbf{加粗}、\textit{斜体}、\underline{下划线}}；\par
{\fangsong 仿宋测试：\textbf{加粗}、\textit{斜体}、\underline{下划线}}。

\section{三级无序列表测试}
\begin{itemize}
\item 一级测试项 A
\item 一级测试项 B
    \begin{itemize}
    \item 二级测试项 B1
    \item 二级测试项 B2
        \begin{itemize}
        \item 三级测试项 B2-1
        \item 三级测试项 B2-2
            \begin{itemize}
            \item 四级测试项 B2-1a
            \item 四级测试项 B2-2b
            \end{itemize}
        \end{itemize}
    \end{itemize}
\end{itemize}

\section{三级有序列表测试}
\begin{enumerate}
\item 一级有序测试 1
\item 一级有序测试 2
    \begin{enumerate}
    \item 二级有序测试 2-1
    \item 二级有序测试 2-2
        \begin{enumerate}
        \item 三级有序测试 2-2-1
        \item 三级有序测试 2-2-2
            \begin{enumerate}
            \item 四级有序测试 2-2-2-1
            \item 四级有序测试 2-2-2-2
            \end{enumerate}
        \end{enumerate}
    \end{enumerate}
\end{enumerate}

% ========== 第2章：复杂数学公式测试 ==========
\chapter{复杂数学公式测试}
\begin{introduction}
\item 单行复杂公式
\item 多行公式（align）
\item 矩阵/行列式
\item 多重积分/求和/极限/分式
\item 公式标签与交叉引用
\end{introduction}

\section{单行复杂公式}
行内复杂公式：$\sum_{i=1}^n \frac{x_i^2}{\sqrt{y_i+1}} \geq \bar{x}^2$、$\lim_{n\to\infty} \left(1+\frac{1}{n}\right)^n = e$。

\section{多行公式与矩阵}
\begin{align}
a^2 + b^2 &= c^2 \label{eq:gougu} \\
\int_{0}^{1}\int_{0}^{1} f(x,y) dxdy &= 1 \label{eq:double} \\
\begin{vmatrix}
a & b \\
c & d
\end{vmatrix} &= ad-bc \label{eq:det}
\end{align}

\section{大型算子公式}
\begin{equation}\label{eq:series}
\sum_{k=0}^{\infty} \frac{x^k}{k!} = e^x,\quad \oint_{C} Pdx+Qdy = \iint_{D}\left(\frac{\partial Q}{\partial x}-\frac{\partial P}{\partial y}\right)dxdy
\end{equation}

公式交叉引用：勾股定理\eqref{eq:gougu}、二重积分\eqref{eq:double}、行列式\eqref{eq:det}、级数\eqref{eq:series}。

% ========== 第3章：全部数学环境测试 ==========
\chapter{全数学环境测试}
\begin{introduction}
\item 定义/定理/引理/推论/公理/公设
\item 命题/示例/问题/习题
\item 提示/备注/结论/假设/性质
\item 解答/证明环境
\end{introduction}

\section{环境测试}

% 1. 定义
\ifdefined\simpleTest
\begin{definition}[测试定义]\label{def:test}
\else
\begin{definition}[测试定义]{test}
\fi
这是定义环境测试内容，无实际数学意义，仅验证样式。
\end{definition}

% 2. 定理
\ifdefined\simpleTest
\begin{theorem}[测试定理]\label{thm:test}
\else
\begin{theorem}[测试定理]{test}
\fi
若 $x>0$，则 $x+1>1$，对应公式\eqref{eq:gougu}。
\end{theorem}

% 3. 引理
\ifdefined\simpleTest
\begin{lemma}[测试引理]\label{lem:test}
\else
\begin{lemma}[测试引理]{test}
\fi
对任意实数 $x$，有 $x^2\geq0$。
\end{lemma}

% 4. 推论
\ifdefined\simpleTest
\begin{corollary}[测试推论]\label{cor:test}
\else
\begin{corollary}[测试推论]{test}
\fi
由引理\ref{lem:test}可得：$(x-1)^2\geq0$。
\end{corollary}

% 5. 公理
\ifdefined\simpleTest
\begin{axiom}[测试公理]\label{axi:test}
\else
\begin{axiom}[测试公理]{test}
\fi
测试用公理，无实际含义。
\end{axiom}

% 6. 公设
\ifdefined\simpleTest
\begin{postulate}[测试公设]\label{pos:test}
\else
\begin{postulate}[测试公设]{test}
\fi
测试用公设，仅验证环境。
\end{postulate}

% 7. 命题
\ifdefined\simpleTest
\begin{proposition}[测试命题]\label{pro:test}
\else
\begin{proposition}[测试命题]{test}
\fi
测试命题内容，关联定义\ref{def:test}。
\end{proposition}

% 8. 示例
\begin{example}\label{exa:test}
示例环境测试，自动编号。
\end{example}

% 9. 问题
\begin{problem}\label{probl:test}
问题环境测试，用于提出待解决问题。
\end{problem}

% 10. 习题
\begin{exercise}\label{exer:test}
计算：$1+2+3+\dots+10$。
\end{exercise}

% 11. 提示
\begin{note}
提示环境：无编号，用于关键提醒。
\end{note}

% 12. 备注
\begin{remark}
备注环境：补充说明用。
\end{remark}

% 13. 结论
\begin{conclusion}
结论环境：测试总结性内容。
\end{conclusion}

% 14. 假设
\begin{assumption}
假设环境：测试条件设定。
\end{assumption}

% 15. 性质
\begin{property}
性质环境：测试对象特性。
\end{property}

% 16. 解答
\begin{solution}
习题\ref{exer:test}解答：和为 $55$（测试答案）。
\end{solution}

% 17. 证明
\begin{proof}
定理\ref{thm:test}证明：因 $x>0$，两边加1即得证。
\end{proof}

% ========== 第4章：交叉引用·图表·旁注测试 ==========
\chapter{交叉引用与拓展功能}
\begin{introduction}
\item 全类型交叉引用
\item 简单图表测试
\item 旁注功能测试
\item 章节跳转验证
\end{introduction}

\section{全交叉引用汇总}
定义\ref{def:test}、定理\ref{thm:test}、引理\ref{lem:test}、推论\ref{cor:test}、公理\ref{axi:test}、公设\ref{pos:test}、命题\ref{pro:test}、示例\ref{exa:test}、问题\ref{probl:test}、习题\ref{exer:test}。

\section{图表测试}
\begin{figure}[h]
    \centering
    \includegraphics[width=0.5\textwidth]{../../image/scatter.jpg}
    \caption{测试图}\label{fig:test}
\end{figure}
\begin{table}[h]
    \centering
    \begin{tabular}{c*{6}{c}}
    \toprule
        $n$ & 1 & 2 & 3 & 4 & 5 & 6 \\
    \midrule
        $a_n$ & 1 & 1 & 2 & 3 & 5 & 8 \\ 
    \bottomrule
    \end{tabular}
    \caption{测试表}\label{tab:test}
\end{table}
图表引用：\figref{fig:test}为测试图片，\tabref{tab:test}为测试表。

% ========== 章后习题 ==========
\begin{problemset}[本章测试习题]
\item 验证定义\ref{def:test}的显示样式
\item 推导公式\eqref{eq:double}
\item 查看图\ref{fig:test}是否正常显示
\end{problemset}

% ========== 参考文献 ==========
\nocite{*}
\printbibliography

% ========== 附录 ==========
\appendix
\chapter{测试附录}
附录内容用于验证附录章节、页眉、页码、排版样式，无实际信息。
\section{附录测试一}
附录内容用于验证附录章节、页眉、页码、排版样式，无实际信息。
\newpage
\section{附录测试二}
附录内容用于验证附录章节、页眉、页码、排版样式，无实际信息。
\end{document}